% Emacs, this is -*-latex-*-

\ex{Estimativa de m\'axima verossimilhan\'ca para uma Gaussiana}
\label{ex:Gauss-MLE}

A fdp Gaussiana parametrizada pela m\'edia $\mu$ e desvio padr\~ao $\sigma$ \'e dada por
$$ p(x; \thetab) = \frac{1}{\sqrt{2\pi\sigma^2}}\exp\left[-\frac{(x-\mu)^2}{2\sigma^2}\right], \quad \quad \thetab = (\mu, \sigma).$$

\begin{exenumerate}
    
    \item Dados dados iid $\data = \{x_1, \ldots, x_n\}$, qual \'e a fun\'c\~ao de verossimilhan\'ca $L(\thetab)$ para o modelo Gaussiano?
    
    \begin{solution}
        Para dados iid, a fun\'c\~ao de verossimilhan\'ca \'e
        \begin{align}
            L(\thetab) &= \prod_i^n p(x_i;\thetab) \\
            &= \prod_i^n \frac{1}{\sqrt{2\pi\sigma^2}}\exp\left[-\frac{(x_i-\mu)^2}{2\sigma^2}\right]\\
            &= \frac{1}{(2\pi\sigma^2)^{n/2}}\exp\left[-\frac{1}{2\sigma^2}\sum_{i=1}^n (x_i-\mu)^2\right].
        \end{align}
        
    \end{solution}
    
    \item Qual \'e a fun\'c\~ao de log-verossimilhan\'ca $\ell(\thetab)$?
    
    \begin{solution}
        
        Tomando o logaritmo da fun\'c\~ao de verossimilhan\'ca, obtemos
        \begin{align}
            \ell(\thetab) & = -\frac{n}{2} \log(2\pi \sigma^2) -\frac{1}{2\sigma^2} \sum_{i=1}^n (x_i-\mu)^2
        \end{align}
        
    \end{solution}
    
    \item Mostre que as estimativas de m\'axima verossimilhan\'ca para a m\'edia $\mu$ e o desvio padr\~ao $\sigma$ s\~ao a m\'edia amostral 
    \begin{equation}
        \bar{x}= \frac{1}{n} \sum_{i=1}^n x_i
    \end{equation}
    e a raiz quadrada da vari\^ancia amostral
    \begin{equation}
        S^2 = \frac{1}{n} \sum_{i=1}^n (x_i -\bar{x})^2.
    \end{equation}
    
    \begin{solution}
        Como o logaritmo \'e estritamente monotonicamente crescente, o
        maximizador da log-verossimilhan\'ca \'e igual ao maximizador da
        verossimilhan\'ca. \'E mais f\'acil tirar derivadas para a
        fun\'c\~ao de log-verossimilhan\'ca do que para a fun\'c\~ao de verossimilhan\'ca, de modo que
        a estimativa de m\'axima verossimilhan\'ca \'e tipicamente determinada usando a log-verossimilhan\'ca.
        
        Dada a express\~ao alg\'ebrica de $\ell(\thetab)$, \'e mais simples
        trabalhar com a vari\^ancia $v=\sigma^2$ em vez do desvio
        padr\~ao. Como $\sigma > 0$, a fun\'c\~ao $v = g(\sigma) = \sigma^2$
        \'e invers\'ivel, e a invari\^ancia do MLE \`a reparametriza\'c\~ao
        garante que
        $$ \hat{\sigma} = \sqrt{\hat{v}}. $$
        
        Agora, portanto, maximizamos a fun\'c\~ao $J(\mu, v)$,
        \begin{align}
            J(\mu, v) & = -\frac{n}{2} \log(2\pi v) -\frac{1}{2v} \sum_{i=1}^n (x_i-\mu)^2
        \end{align}
        em rela\'c\~ao a $\mu$ e $v$.
        
        Tomando as derivadas parciais, obtemos
        \begin{align}
            \frac{\partial J }{\partial\mu} & = \frac{1}{v} \sum_{i=1}^n (x_i-\mu)\\
            &= \frac{1}{v} \sum_{i=1}^n x_i  -\frac{n}{v} \mu\\
            \frac{\partial J}{\partial v} &= -\frac{n}{2} \frac{1}{v} + \frac{1}{2v^2} \sum_{i=1}^n (x_i-\mu)^2
        \end{align}
        Uma condi\'c\~ao necess\'aria para a otimalidade \'e que as derivadas parciais sejam zero. Obtemos assim as condi\'c\~oes
        \begin{align}
            \frac{1}{v} \sum_{i=1}^n (x_i-\mu) &=0\\
            -\frac{n}{2} \frac{1}{v} + \frac{1}{2v^2} \sum_{i=1}^n (x_i-\mu)^2 & = 0
        \end{align}
        Da primeira condi\'c\~ao segue que
        \begin{align}
            \hat{\mu} &= \frac{1}{n}\sum_{i=1}^n x_i
        \end{align}
        A segunda condi\'c\~ao torna-se, ent\~ao,
        \begin{align}
            -\frac{n}{2} \frac{1}{v} + \frac{1}{2v^2} \sum_{i=1}^n (x_i-\hat{\mu})^2 & = 0 \quad \quad (\text{multiplique por } v^2 \text{ e reorganize})\\
            \frac{1}{2} \sum_{i=1}^n (x_i-\hat{\mu})^2 & = \frac{n}{2} v,
        \end{align}
        e, portanto,
        \begin{align}
            \hat{v} & = \frac{1}{n} \sum_{i=1}^n (x_i-\hat{\mu})^2,
        \end{align}
        Verificamos agora que esta solu\'c\~ao corresponde a um m\'aximo calculando a matriz Hessiana
        \begin{align}
            \H(\mu, v) & = \begin{pmatrix}
                \frac{\partial^2 J }{\partial \mu^2} & \frac{\partial^2 J }{\partial \mu \partial v}\\
                \frac{\partial^2 J }{\partial \mu \partial v} & \frac{\partial^2 J }{\partial v^2}
            \end{pmatrix}
        \end{align}
        Se a Hessiana for definida negativa em $(\hat{\mu}, \hat{v})$, o
        ponto \'e um m\'aximo (local). Como temos apenas um ponto cr\'itico,
        $(\hat{\mu}, \hat{v})$, o m\'aximo local \'e tamb\'em um
        m\'aximo global. Tomando as segundas derivadas, obtemos
        \begin{align}
            \H(\mu, v) & = \begin{pmatrix}
                -\frac{n}{v} & -\frac{1}{v^2} \sum_{i=1}^n (x_i-\mu)\\
                -\frac{1}{v^2} \sum_{i=1}^n (x_i-\mu) & \frac{n}{2} \frac{1}{v^2} - \frac{1}{v^3} \sum_{i=1}^n (x_i-\mu)^2
            \end{pmatrix}.
        \end{align}
        Substituindo os valores para $(\hat{\mu}, \hat{v})$ resulta em
        \begin{align}
            \H(\hat{\mu}, \hat{v}) & = \begin{pmatrix}
                -\frac{n}{\hat{v}} & 0 \\
                0 & -\frac{n}{2} \frac{1}{\hat{v}^2}
            \end{pmatrix},
        \end{align}
        que \'e definida negativa. Note que a curvatura (negativa)
        aumenta com $n$, o que significa que $J(\mu, v)$, e
        consequentemente a log-verossimilhan\'ca, torna-se cada vez mais pontiaguda \`a medida que o
        n\'umero de pontos de dados $n$ aumenta.
    \end{solution}
    
\end{exenumerate}



\ex{Posterior da m\'edia de uma Gaussiana com vari\^ancia conhecida} Dados dados iid $\data = \{x_1, \ldots, x_n\}$, compute $p(\mu | \data, \sigma^2)$ para o
modelo Bayesiano
\begin{align}
    p(x | \mu) &= \frac{1}{\sqrt{2\pi\sigma^2}}\exp\left[-\frac{(x-\mu)^2}{2\sigma^2}\right] &
    p(\mu; \mu_0, \sigma_0^2) &=  \frac{1}{\sqrt{2\pi\sigma_0^2}}\exp\left[-\frac{(\mu-\mu_0)^2}{2\sigma_0^2}\right]
\end{align}
onde $\sigma^2$ \'e uma quantidade fixa conhecida.\\
Dica: Voc\^e pode usar que
\begin{equation}
    \Gauss(x ; m_1, \sigma_1^2) \Gauss(x ; m_2, \sigma_2^2) \propto \Gauss(x ; m_3, \sigma_3^2) \label{eq:Gaussian-mult-model-based-learning}
\end{equation}
onde 
\begin{align}
    \Gauss(x ; \mu, \sigma^2) &= \frac{1}{\sqrt{2 \pi \sigma^2}} \exp\left[ -\frac{(x-\mu)^2}{2\sigma^2} \right]\\
    \sigma^2_3 & = \left( \frac{1}{\sigma_1^2} + \frac{1}{\sigma_2^2}\right)^{-1} = \frac{\sigma_1^2 \sigma_2^2}{\sigma_1^2+\sigma_2^2}\\
    m_3 & = \sigma_3^2\left( \frac{m_1}{\sigma_1^2}+ \frac{m_2}{\sigma_2^2}\right) = m_1 + \frac{\sigma_1^2}{\sigma_1^2+\sigma_2^2}(m_2-m_1) \label{eq:mu-mult-alt}
\end{align}

\begin{solution}
    Reutilizamos a express\~ao para a verossimilhan\'ca $L(\mu)$ do \exref{ex:Gauss-MLE}.
    \begin{equation}
        L(\mu) = \frac{1}{(2\pi\sigma^2)^{n/2}}\exp\left[-\frac{1}{2\sigma^2}\sum_{i=1}^n (x_i-\mu)^2\right],
    \end{equation}
    que podemos escrever como
    \begin{align}
        L(\mu) & \propto \exp\left[-\frac{1}{2\sigma^2}\sum_{i=1}^n (x_i-\mu)^2\right]\\
        & \propto  \exp\left[-\frac{1}{2\sigma^2}\sum_{i=1}^n (x_i^2-2\mu x_i + \mu^2)\right]\\
        & \propto  \exp\left[-\frac{1}{2\sigma^2}\left( -2 \mu \sum_{i=1}^n x_i + n \mu^2 \right)\right]\\
        & \propto  \exp\left[-\frac{1}{2\sigma^2}\left( -2n \mu \xBar + n \mu^2 \right)\right]\\
        & \propto  \exp\left[-\frac{n}{2\sigma^2}(\mu-\xBar)^2 \right]\\
        & \propto  \Gauss(\mu ; \xBar, \sigma^2/n).
    \end{align}
    A posterior \'e
    \begin{align}
        p(\mu | \data) & \propto L(\theta) p(\mu; \mu_0, \sigma_0^2)\\
        & \propto  \Gauss(\mu ; \xBar, \sigma^2/n)  \Gauss(\mu ; \mu_0, \sigma_0^2)
    \end{align}
    de modo que, com \eqref{eq:Gaussian-mult-model-based-learning}, temos
    \begin{align}
        p(\mu | \data) & \propto \Gauss(\mu ; \mu_n, \sigma_n^2)\\
        \sigma^2_n & =  \left( \frac{1}{\sigma^2/n} + \frac{1}{\sigma_0^2}\right)^{-1} \\
        & =  \frac{ \sigma_0^2 \sigma^2/n}{\sigma_0^2+\sigma^2/n}\\
        \mu_n & =  \sigma_n^2\left( \frac{\xBar}{\sigma^2/n}+ \frac{\mu_0}{\sigma_0^2}\right)\\
        & = \frac{1}{\sigma_0^2+\sigma^2/n}\left(\sigma_0^2 \xBar + (\sigma^2/n) \mu_0\right)\\
        & = \frac{\sigma_0^2}{\sigma_0^2+\sigma^2/n} \xBar + \frac{\sigma^2/n} {\sigma_0^2+\sigma^2/n} \mu_0.
    \end{align}
    \`A medida que $n$ aumenta, $\sigma^2/n$ tende a zero, de modo que $\sigma_n^2 \to 0$
    e $\mu_n \to \xBar$. Isso significa que, com uma quantidade crescente de
    dados, a posterior da m\'edia tende a concentrar-se em torno da
    estimativa de m\'axima verossimilhan\'ca $\xBar$.
    
    De \eqref{eq:mu-mult-alt}, temos tamb\'em que
    \begin{align}
        \mu_n & = \mu_0 + \frac{\sigma_0^2}{\sigma^2/n+\sigma_0^2}(\xBar-\mu_0),
    \end{align}
    o que mostra mais claramente que o valor de $\mu_n$ encontra-se em uma linha com extremidades $\mu_0$ (para $n = 0$) e $\xBar$ (para $n \to \infty$). \`A medida que
    a quantidade de dados aumenta, $\mu_n$ move-se da m\'edia sob a priori, $\mu_0$, para a m\'edia da amostra observada, que \'e o MLE $\xBar$.
\end{solution}


% --------------------------------------------------------------- 

\ex{Estimativa de m\'axima verossimilhan\'ca de tabelas de probabilidade em modelos gr\'aficos direcionados totalmente observados de vari\'aveis bin\'arias}
\label{ex:MLE-DGM-alt}
Assumimos que nos \'e dado um modelo gr\'afico direcionado parametrizado para vari\'aveis
$x_1, \ldots, x_d$, 
\begin{align}
    p(\x; \thetab) &= \prod_{i=1}^d p(x_i | \pa_i; \thetab_i) \quad \quad x_i \in \{0,1\} \label{eq:p-dag-alt}
\end{align}
onde os condicionais s\~ao representados por tabelas de probabilidade parametrizadas. Por exemplo, se $\pa_3 = \{x_1, x_2\}$, $p(x_3 | \pa_3; \thetab_3)$ \'e representado como

\begin{center}
    \begin{tabular}{@{}cll@{}}
        \toprule
        $p(x_3=1 | x_1, x_2; \theta^1_3, \ldots, \theta_3^4))$ & $x_1$ & $x_2$\\
        \midrule
        $\theta^1_3$ & 0 & 0\\
        $\theta^2_3$ & 1 & 0\\
        $\theta^3_3$ & 0 & 1\\
        $\theta^4_3$ & 1 & 1\\
        \bottomrule
    \end{tabular}
\end{center}
com $\thetab_3 = (\theta_3^1,\theta_3^2, \theta_3^3, \theta_3^4)$,
e onde os sobrescritos $j$ de $\theta_3^j$ enumeram os diferentes estados
em que os pais podem estar.

\begin{exenumerate}
    
    \item Assumindo que $x_i$ possua $m_i$ pais, verifique que a parametriza\'c\~ao por tabela de $p(x_i | \pa_i; \thetab_i)$ \'e equivalente a escrever $p(x_i | \pa_i; \thetab_i)$ como
    \begin{align}
        p(x_i | \pa_i; \thetab_i) &= \prod_{s=1}^{S_i} (\theta_i^s)^{\ind(x_i=1, \pa_i=s)} (1-\theta_i^s)^{\ind(x_i=0, \pa_i=s)} \label{eq:pcond-fac-alt}
    \end{align}
    onde $S_i = 2^{m_i}$ \'e o n\'umero total de estados/configura\'c\~oes em que os
    pais podem estar, e $\ind(x_i=1, \pa_i=s)$ \'e um se $x_i=1$ e $\pa_i=s$, e zero caso contr\'ario.
    
    \begin{solution}
        O n\'umero de configura\'c\~oes em que $m$ pais bin\'arios podem estar \'e dado por $S_i$. A quest\~ao, portanto, resume-se a mostrar que $p(x_i=1 | \pa_i=k; \thetab_i) = \theta_i^k$ para qualquer estado $k \in \{1, \ldots, S_i\}$ dos pais de $x_i$. Como $\ind(x_i=1, \pa_i=s)=0$ exceto se $s=k$, temos de fato que
        \begin{align}
            p(x_i = 1| \pa_i = k; \thetab_i) &= \left[\prod_{s \neq k} (\theta_i^s)^0 (1-\theta_i^s)^0\right] (\theta_i^k)^{\ind(x_i=1, \pa_i=k)} (1-\theta_i^k)^{\ind(x_i=0, \pa_i=k)}\\
            &= 1 \cdot  (\theta_i^k)^{\ind(x_i=1, \pa_i=k)} (1-\theta_i^k)^0\\
            & = \theta_i^k.
        \end{align}
        
    \end{solution}
    
    \item Para dados iid $\data = \{\x^{(1)}, \ldots, \x^{(n)}\}$, mostre que a verossimilhan\'ca pode ser representada como
    \begin{align}
        p(\data; \thetab) & = \prod_{i=1}^d \prod_{s=1}^{S_i}  (\theta_i^s)^{n_{x_i=1}^s} (1-\theta_i^s)^{n_{x_i=0}^s} \label{eq:joint-dag-alt}
    \end{align}
    onde $n_{x_i=1}^s$ \'e o n\'umero de vezes que o padr\~ao $(x_i=1, \pa_i=s)$ ocorre nos dados $\data$, e equivalentemente para $n_{x_i=0}^s$.
    
    \begin{solution}
        Como os dados s\~ao iid, temos
        \begin{align}
            p(\data; \thetab) & = \prod_{j=1}^n p(\x^{(j)}; \thetab) \\
        \end{align}
        onde cada termo $p(\x^{(j)}; \thetab)$ se fatoriza como em \eqref{eq:p-dag-alt},
        \begin{align}
            p(\x^{(j)}; \thetab) &= \prod_{i=1}^d p(x_i^{(j)} | \pa_i^{(j)}; \thetab_i)
        \end{align}
        com $x_i^{(j)}$ denotando o $i$-\'esimo elemento de $\x^{(j)}$ e $\pa_i^{(j)}$ os pais correspondentes. Os condicionais $p(x_i^{(j)} | \pa_i^{(j)}; \thetab_i)$ fatorizam-se ainda de acordo com \eqref{eq:pcond-fac-alt},
        \begin{align}
            p(x_i^{(j)} | \pa_i^{(j)}; \thetab_i) &= \prod_{s=1}^{S_i} (\theta_i^s)^{\ind(x_i^{(j)}=1, \pa_i^{(j)}=s)} (1-\theta_i^s)^{\ind(x_i^{(j)}=0, \pa_i^{(j)}=s)},
        \end{align}
        de modo que
        \begin{align}
            p(\data; \thetab) & = \prod_{j=1}^n \prod_{i=1}^d p(x_i^{(j)} | \pa_i^{(j)}; \thetab_i)\\
            & =  \prod_{j=1}^n  \prod_{i=1}^d \prod_{s=1}^{S_i} (\theta_i^s)^{\ind(x_i^{(j)}=1, \pa_i^{(j)}=s)} (1-\theta_i^s)^{\ind(x_i^{(j)}=0, \pa_i^{(j)}=s)}
        \end{align}
        Trocando a ordem dos produtos para que o produto sobre os pontos de dados venha primeiro, obtemos
        \begin{align}
            p(\data; \thetab)  & = \prod_{i=1}^d \prod_{s=1}^{S_i} \prod_{j=1}^n(\theta_i^s)^{\ind(x_i^{(j)}=1, \pa_i^{(j)}=s)} (1-\theta_i^s)^{\ind(x_i^{(j)}=0, \pa_i^{(j)}=s)}
        \end{align}
        Dividimos em seguida o produto sobre $j$ em dois produtos, um para todos os $j$ onde $x_i^{(j)}=1$, e outro para todos os $j$ onde $x_i^{(j)}=0$
        \begin{align}
            p(\data; \thetab)   &=  \prod_{i=1}^d \prod_{s=1}^{S_i} \prod_{\substack{j: \\x_i^{(j)}=1}}\prod_{\substack{j: \\x_i^{(j)}=0}} (\theta_i^s)^{\ind(x_i^{(j)}=1, \pa_i^{(j)}=s)} (1-\theta_i^s)^{\ind(x_i^{(j)}=0, \pa_i^{(j)}=s)}\\
            &=  \prod_{i=1}^d \prod_{s=1}^{S_i} \prod_{\substack{j: \\x_i^{(j)}=1}} (\theta_i^s)^{\ind(x_i^{(j)}=1, \pa_i^{(j)}=s)} \prod_{\substack{j: \\x_i^{(j)}=0}}(1-\theta_i^s)^{\ind(x_i^{(j)}=0, \pa_i^{(j)}=s)}\\
            & = \prod_{i=1}^d \prod_{s=1}^{S_i} (\theta_i^s)^{ \sum_{j=1}^n \ind(x_i^{(j)}=1, \pa_i^{(j)}=s)} (1-\theta_i^s)^{ \sum_{j=1}^n\ind(x_i^{(j)}=0, \pa_i^{(j)}=s)}\\
            & =  \prod_{i=1}^d \prod_{s=1}^{S_i}  (\theta_i^s)^{n_{x_i=1}^s} (1-\theta_i^s)^{n_{x_i=0}^s}
        \end{align}
        onde
        \begin{align}
            n_{x_i=1}^s &=  \sum_{j=1}^n\ind(x_i^{(j)}=1, \pa_i^{(j)}=s) & n_{x_i=0}^s &=  \sum_{j=1}^n\ind(x_i^{(j)}=0, \pa_i^{(j)}=s)
        \end{align}
        \'e o n\'umero de vezes que $x_i=1$ e $x_i=0$, respectivamente, com seus pais no estado $s$.
        
    \end{solution}
    
    \item \label{q:loglik-fully-observed-dgm-sum-bernoulli-alt}Mostre que a log-verossimilhan\'ca se decomp\~oe em somas de termos que podem ser independentemente otimizados, e que cada termo corresponde \`a log-verossimilhan\'ca de um modelo de Bernoulli.
    
    \begin{solution}
        A log-verossimilhan\'ca $\ell(\thetab)$ \'e igual a
        \begin{align}
            \ell(\thetab) & = \log p(\data; \thetab) \\
            & = \log \prod_{i=1}^d \prod_{s=1}^{S_i}  (\theta_i^s)^{n_{x_i=1}^s} (1-\theta_i^s)^{n_{x_i=0}^s}\\
            & = \sum_{i=1}^d \sum_{s=1}^{S_i} \log \left[  (\theta_i^s)^{n_{x_i=1}^s} (1-\theta_i^s)^{n_{x_i=0}^s} \right]\\
            & = \sum_{i=1}^d \sum_{s=1}^{S_i} n_{x_i=1}^s \log (\theta_i^s) + n_{x_i=0}^s \log (1-\theta_i^s)
        \end{align}
        Como os par\^ametros $\theta_i^s$ n\~ao est\~ao acoplados de forma alguma, a maximiza\'c\~ao de $\ell(\thetab)$ pode ser alcan\'cada maximizando cada termo $\ell_{is}(\theta_i^s)$ individualmente,
        \begin{align}
            \ell_{is}(\theta_i^s) & = n_{x_i=1}^s \log (\theta_i^s) + n_{x_i=0}^s \log (1-\theta_i^s).
        \end{align}
        Al\'em disso, $\ell_{is}(\theta_i^s)$ corresponde \`a log-verossimilhan\'ca de um modelo de Bernoulli com probabilidade de sucesso $\theta_i^s$ e dados com n\'umero $n_{x_i=1}^s$ de uns e n\'umero $n_{x_i=0}^s$ de zeros.
        
    \end{solution}
    
    \item \label{q:mle-bernoulli-alt} Determine a estimativa de m\'axima verossimilhan\'ca $\hat{\theta}$ para o
    modelo de Bernoulli
    \begin{align}
        p(x; \theta) &= \theta^{x} (1-\theta)^{1-x},& \theta &\in [0,\;1],&  x &\in\{0,1\} 
    \end{align}
    e dados iid $x_1, \ldots, x_n$.
    \begin{solution}
        A fun\'c\~ao de log-verossimilhan\'ca \'e
        \begin{align}      
            \ell(\theta) & = \sum_{i=1}^n \log p(x_i; \theta) \\
            & = \sum_{i=1}^n x_i \log(\theta) + (1-x_i) \log(1-\theta).
        \end{align}
        Como $\log(\theta)$ e $\log(1-\theta)$ n\~ao dependem de $i$, podemos pux\'a-los para fora da soma e a fun\'c\~ao de log-verossimilhan\'ca pode ser escrita como
        \begin{equation}
            \ell(\theta) = n_{x=1} \log (\theta) + n_{x=0} \log(1-\theta)
        \end{equation}
        onde $n_{x=1} = \sum_{i=1}^n x_i = \sum_{i=1}^n \ind(x_i=1)$ e
        $n_{x=0} = n-n_{x=1}$ s\~ao o n\'umero de uns e zeros nos
        dados. Como $\theta \in [0,\;1]$, temos que resolver o problema de
        otimiza\'c\~ao restrito
        \begin{equation}
            \hat{\theta} = \argmax_{\theta \in [0,\;1]}   \ell(\theta)
        \end{equation}
        Existem v\'arias formas de resolver o problema. Uma op\'c\~ao \'e
        determinar o otimizador \emph{irrestrito} e ent\~ao verificar se ele
        satisfaz a restri\'c\~ao. A primeira derivada \'e igual a
        \begin{equation}
            \ell'(\theta) = \frac{n_{x=1}}{\theta} - \frac{n_{x=0}}{1-\theta}
        \end{equation}
        e a segunda derivada \'e
        \begin{equation}
            \ell''(\theta) = -\frac{n_{x=1}}{\theta^2} - \frac{n_{x=0}}{(1-\theta)^2}
        \end{equation}
        A segunda derivada \'e sempre negativa para $\theta \in (0,1)$,
        o que significa que $\ell(\theta)$ \'e estritamente c\^oncava em $(0,1)$
        e que um otimizador que n\~ao esteja na fronteira corresponde a um
        m\'aximo. Definindo a primeira derivada como zero, obtemos a condi\'c\~ao
        \begin{equation}
            \frac{n_{x=1}}{\theta} = \frac{n_{x=0}}{1-\theta}
        \end{equation}
        Resolvendo para $\theta$ resulta em
        \begin{align}
            (1-\theta) n_{x=1} &= n_{x=0} \theta\\
        \end{align}
        de modo que
        \begin{align}
            n_{x=1} & = \theta( n_{x=0}+ n_{x=1})\\
            & = \theta n
        \end{align}
        Assim, encontramos
        \begin{equation}
            \hat{\theta} = \frac{n_{x=1}}{n}.
        \end{equation}
        Para $n_{x=1}< n$, temos $\hat{\theta} \in (0,1)$, de modo que a
        restri\'c\~ao n\~ao est\'a, de fato, ativa.
        
        Na deriva\'c\~ao, tivemos que excluir os casos de fronteira onde $\theta$
        \'e 0 ou 1. Notamos que, por exemplo, $\hat{\theta}=1$ \'e obtido quando
        $n_{x=1}=n$, ou seja, quando observamos apenas 1s no conjunto de dados. Nesse
        caso, $n_{x=0}=0$ e a fun\'c\~ao de log-verossimilhan\'ca \'e igual a
        $n\log(\theta)$, que \'e estritamente crescente e, portanto, atinge o
        m\'aximo em $\hat{\theta}=1$. Um argumento semelhante mostra que se
        $n_{x=1}=0$, o m\'aximo ocorre em $\hat{\theta}=0$. Logo, a estimativa de
        m\'axima verossimilhan\'ca
        \begin{equation}
            \hat{\theta} = \frac{n_{x=1}}{n}
        \end{equation}
        \'e v\'alida para todo $n_{x=1} \in \{0, \ldots, n\}$.
        
        Uma abordagem alternativa para lidar com a restri\'c\~ao \'e
        reparametrizar a fun\'c\~ao objetivo e trabalhar com o logit (log-odds) $\eta$,
        \begin{equation}
            \eta = g(\theta) = \log\left[ \frac{\theta}{1-\theta}\right].
        \end{equation}
        O logit assume valores em $\mathbb{R}$, de modo que $\eta$ \'e
        irrestrito. A transforma\'c\~ao de $\theta$ para $\eta$ \'e
        invers\'ivel e
        \begin{equation}
            \theta = g^{-1}(\eta) = \frac{\exp(\eta)}{1+\exp(\eta)} = \frac{1}{1+\exp(-\eta)}.
        \end{equation}
        O problema de otimiza\'c\~ao torna-se ent\~ao
        \begin{align*}
            \hat{\eta} &= \argmax_{\eta}  \{ n_{x=1} \eta - n \log(1+\exp(\eta)) \}
        \end{align*}
        Calcular a segunda derivada mostra que o objetivo \'e c\^oncavo
        para todo $\eta$ e o maximizador $\hat{\eta}$ pode ser determinado
        definindo a primeira derivada como zero. A estimativa de m\'axima
        verossimilhan\'ca de $\theta$ \'e ent\~ao dada por
        \begin{equation}
            \hat{\theta} = \frac{\exp(\hat{\eta})}{1+\exp(\hat{\eta})}
        \end{equation}
        A raz\~ao para isto \'e a seguinte: Seja $J(\eta) = \ell(
        g^{-1}(\eta))$ a log-verossimilhan\'ca vista como uma fun\'c\~ao de
        $\eta$. Como $g$ e $g^{-1}$ s\~ao invers\'iveis, temos que
        \begin{align}
            \max_{\theta \in [0,1]} \ell(\theta) &= \max_\eta J(\eta) \\
            \argmax_{\theta \in [0,1]} \ell(\theta) &= g^{-1}\left( \argmax_{\eta} J(\eta) \right).
        \end{align}
    \end{solution}
    
    \item Retornando ao modelo gr\'afico direcionado totalmente observado, conclua que as estimativas de m\'axima verossimilhan\'ca s\~ao dadas por
    \begin{align}
        \hat{\theta}_i^s &= \frac{n_{x_i=1}^s}{n_{x_i=1}^s+n_{x_i=0}^s} = \frac{ \sum_{j=1}^n\ind(x_i^{(j)}=1, \pa_i^{(j)}=s) }{ \sum_{j=1}^n\ind(\pa_i^{(j)}=s) }
    \end{align}
    
    \begin{solution}
        Dado o resultado da quest\~ao
        \ref{q:loglik-fully-observed-dgm-sum-bernoulli-alt}, podemos otimizar
        cada termo $\ell_{is}(\theta_i^s)$ separadamente. Cada termo
        corresponde formalmente a uma log-verossimilhan\'ca de um modelo de Bernoulli, de modo que
        podemos usar os resultados da quest\~ao \ref{q:mle-bernoulli-alt} para obter
        \begin{align}
            \hat{\theta}_i^s &= \frac{n_{x_i=1}^s}{n_{x_i=1}^s+n_{x_i=0}^s}.
        \end{align}
        Como $n_{x_i=1}^s = \sum_{j=1}^n\ind(x_i^{(j)}=1, \pa_i^{(j)}=s)$ e
        \begin{align}
            n_{x_i=1}^s+n_{x_i=0}^s &=  \sum_{j=1}^n\ind(x_i^{(j)}=1, \pa_i^{(j)}=s) +  \sum_{j=1}^n\ind(x_i^{(j)}=0, \pa_i^{(j)}=s)\\
            & =  \sum_{j=1}^n\ind( \pa_i^{(j)}=s),
        \end{align}
        temos ainda
        \begin{align}
            \hat{\theta}_i^s &= \frac{ \sum_{j=1}^n\ind(x_i^{(j)}=1, \pa_i^{(j)}=s) }{ \sum_{j=1}^n\ind(\pa_i^{(j)}=s) }.
            \label{eq:dag-table-mle-alt}
        \end{align}
        Portanto, para determinar $\hat{\theta}_i^s$, primeiro contamos o
        n\'umero de vezes que os pais de $x_i$ est\~ao no estado $s$, o que
        fornece o denominador, e ent\~ao, entre estes, contamos o n\'umero de
        vezes que $x_i=1$, o que fornece o numerador.
        
    \end{solution}
    
\end{exenumerate}

% --------------------------------------------------

\ex{Exemplo c\^ancer-amianto-fumo: MLE}
\label{ex:cancer-smoking-asbestos-mle-alt} 
Considere o modelo especificado pelo DAG

\begin{center}
    \begin{tikzpicture}[dgraph]
        \node[cont] (x) at (0,0) {a};
        \node[cont] (y) at (2,0) {s};
        \node[cont] (z) at (1,-1) {c};
        \draw(x) -- (z);
        \draw(y) -- (z);
    \end{tikzpicture}
\end{center}

As distribui\'c\~oes de $a$ e $s$ s\~ao distribui\'c\~oes de Bernoulli com
par\^ametro (probabilidade de sucesso) $\theta_a$ e $\theta_s$,
respectivamente, i.e.\
\begin{equation}
    p(a; \theta_a) = \theta_a^a(1-\theta_a)^{1-a} \quad \quad p(s; \theta_s) = \theta_s^s(1-\theta_s)^{1-s},
\end{equation}
e a distribui\'c\~ao de $c$ dados os pais \'e parametrizada conforme especificado na tabela seguinte
\begin{center}
    \begin{tabular}{@{}cll@{}}
        \toprule
        $p(c=1 | a,s; \theta^1_c, \ldots, \theta_c^4))$ & $a$ & $s$\\
        \midrule
        $\theta^1_c$ & 0 & 0\\
        $\theta^2_c$ & 1 & 0\\
        $\theta^3_c$ & 0 & 1\\
        $\theta^4_c$ & 1 & 1\\
        \bottomrule
    \end{tabular}
\end{center}
Os par\^ametros livres do modelo s\~ao $(\theta_a, \theta_s, \theta^1_c, \ldots, \theta^4_c)$.

Assuma que observamos os seguintes dados iid (cada linha \'e um ponto de dado).

\begin{center}
    \scalebox{1}{
        \begin{tabular}{lll}
            \toprule
            a & s & c\\
            \midrule
            0 &   1 &   1\\
            0 &   0 &   0\\
            1 &   0 &   1\\
            0 &   0 &   0\\
            0 &   1 &   0\\
            \bottomrule
    \end{tabular}}
\end{center}

\begin{exenumerate}
    \item Determine as estimativas de m\'axima verossimilhan\'ca de $\theta_a$ e $\theta_s$
    \begin{solution}
        A estimativa de m\'axima verossimilhan\'ca (MLE) $\hat{\theta}_a$ \'e dada
        pela fra\'c\~ao de vezes que $a$ \'e 1 no conjunto de dados. Portanto,
        $\hat{\theta}_a = 1/5$. Da mesma forma, o MLE $\hat{\theta}_s$ \'e
        $2/5$.
        
    \end{solution}
    
    \item \label{q:cancer-smoking-asbestos-mle-alt-q} Determine as estimativas de m\'axima verossimilhan\'ca de $\theta_c^1, \ldots, \theta_c^4$.
    
    \begin{solution}
        Com \eqref{eq:dag-table-mle-alt}, temos
        \begin{center}
            \begin{tabular}{@{}lll@{}}
                \toprule
                $\hat{p}(c=1 | a,s)$ & $a$ & $s$\\
                \midrule
                $\hat{\theta}^1_c = 0$ & 0 & 0\\
                $\hat{\theta}^2_c = 1/1$ & 1 & 0\\
                $\hat{\theta}^3_c = 1/2$ & 0 & 1\\
                $\hat{\theta}^4_c$ n\~ao definido & 1 & 1\\
                \bottomrule
            \end{tabular}
        \end{center}
        Isso ocorre porque, por exemplo, temos duas observa\'c\~oes onde
        $(a,s)=(0,0)$, e entre elas, $c=1$ nunca ocorre, de modo que o MLE para
        $p(c=1 | a,s)$ \'e zero.
        
        Este exemplo ilustra alguns problemas com estimativas de m\'axima
        verossimilhan\'ca: podemos obter probabilidades extremas, zero ou um, ou, se a
        configura\'c\~ao dos pais n\~ao ocorrer nos dados observados, a estimativa
        ficar\'a indefinida.
        
    \end{solution}
    
\end{exenumerate}

% --------------------------------------------------

\ex{Infer\^encia Bayesiana para o modelo de Bernoulli}
\label{ex:Bayesian-inference-Bernoulli-alt}
Considere o modelo Bayesiano
\begin{align*}
    p(x | \theta) &= \theta^x (1-\theta)^{1-x} & p(\theta ; \alphab_0) = \BetaDist(\theta ; \alpha_0, \beta_0)
\end{align*}
onde $x \in \{0,1\}, \ \theta \in [0,1], \alphab_0 = (\alpha_0,\beta_0)$, e
\begin{equation}
    \BetaDist(\theta; \alpha, \beta) \propto \theta^{\alpha-1}(1-\theta)^{\beta-1} \quad \quad \theta \in [0,1]
\end{equation}

\begin{exenumerate}
    \item Dados dados iid $\data = \{x_1, \ldots, x_n\}$, mostre que a posterior de $\theta$ dado $\data$ \'e
    \begin{align*}
        p(\theta | \data) &= \BetaDist(\theta ; \alpha_n, \beta_n)\\
        \alpha_n &= \alpha_0 + n_{x=1} & \beta_n &= \beta_0 + n_{x=0}
    \end{align*}
    onde $n_{x=1}$ denota o n\'umero de uns e $n_{x=0}$ o n\'umero de zeros nos dados.
    
    \begin{solution}
        Isso decorre de
        \begin{align}
            p(\theta | \data) \propto L(\theta) p(\theta ; \alphab_0) \label{eq:Bernoulli-posterior-def-alt}
        \end{align}
        e da express\~ao para a fun\'c\~ao de verossimilhan\'ca do modelo de Bernoulli, que \'e
        \begin{align}
            L(\theta) &= \prod_{i=1}^n p(x_i|\theta)\\
            & = \prod_{i=1}^n  \theta^{x_i} (1-\theta)^{1-x_i}\\
            & = \theta^{\sum_{i=1}^n x_i} (1-\theta)^{\sum_{i=1}^n (1-x_i)}\\
            & = \theta^{n_{x=1}}(1-\theta)^{n_{x=0}},
        \end{align}
        onde $n_{x=1} = \sum_{i=1}^n x_i$ denota o n\'umero de 1s nos
        dados, e $n_{x=0}=\sum_{i=1}^n (1-x_i)=n-n_{x=1}$ o n\'umero
        de 0s.
        
        Inserir as express\~oes para a verossimilhan\'ca e a priori em
        \eqref{eq:Bernoulli-posterior-def-alt} fornece
        \begin{align}
            p(\theta | \data) &\propto \theta^{n_{x=1}}(1-\theta)^{n_{x=0}}\theta^{\alpha_0-1}(1-\theta)^{\beta_0-1}\\
            & \propto \theta^{\alpha_0+n_{x=1}-1}(1-\theta)^{\beta_0+n_{x=0}-1}\\
            & \propto \BetaDist(\theta, \alpha_0+n_{x=1}, \beta_0+n_{x=0}),
        \end{align}
        que \'e o resultado desejado. Como $\alpha_0$ e $\beta_0$ s\~ao
        atualizados pelas contagens de uns e zeros nos dados, esses
        hiperpar\^ametros s\~ao tamb\'em referidos como
        ``pseudo-contagens''. Alternativamente, pode-se pensar que s\~ao as
        contagens observadas em outro conjunto de dados iid que foi
        previamente analisado e usado para determinar a priori.
    \end{solution}
    
    
    
    \item \label{q:mean-beta-var-alt} Calcule a m\'edia de uma vari\'avel aleat\'oria Beta $f$,
    \begin{equation}
        p(f; \alpha, \beta) =  \BetaDist(f ; \alpha, \beta) \quad \quad f \in [0,1],
    \end{equation}
    usando que
    \begin{equation}
        \int_0^1 f^{\alpha-1} (1-f)^{\beta-1} \ud f = B(\alpha, \beta) = \frac{\Gamma(\alpha)\Gamma(\beta)}{\Gamma(\alpha+\beta)}
    \end{equation}
    onde $ B(\alpha, \beta)$ denota a fun\'c\~ao Beta e onde a fun\'c\~ao Gamma $\Gamma(t)$ \'e definida como
    \begin{equation}
        \Gamma(t) = \int_o^\infty f^{t-1} \exp(-f) \ud f
    \end{equation}
    e satisfaz $\Gamma(t+1) = t \Gamma(t)$.\\
    \emph{Dica: Ser\'a \'util representar a fun\'c\~ao de parti\'c\~ao em termos da fun\'c\~ao Beta.}
    
    \begin{solution}
        Primeiro escrevemos a fun\'c\~ao de parti\'c\~ao de $p(f; \alpha, \beta)$ em termos da fun\'c\~ao Beta
        \begin{align}
            Z(\alpha, \beta) &= \int_0^1 f^{\alpha-1} (1-f)^{\beta-1}\\
            & =  B(\alpha, \beta).
        \end{align}
        Temos ent\~ao que a m\'edia $\E[f]$ \'e dada por
        \begin{align}
            \E[f] & = \int_0^1 f p(f; \alpha, \beta) \ud f \\
            & = \frac{1}{B(\alpha, \beta)} \int_0^1 f  f^{\alpha-1} (1-f)^{\beta-1} \ud f\\
            & = \frac{1}{B(\alpha, \beta)} \int_0^1 f^{\alpha+1-1} (1-f)^{\beta-1} \ud f\\
            & = \frac{B(\alpha+1, \beta)}{B(\alpha, \beta)} \\
            & = \frac{\Gamma(\alpha+1)\Gamma(\beta)}{\Gamma(\alpha+1+\beta)} \frac{\Gamma(\alpha+\beta)}{\Gamma(\alpha)\Gamma(\beta)}\\
            & = \frac{\alpha \Gamma(\alpha)\Gamma(\beta)}{(\alpha+\beta) \Gamma(\alpha+\beta)} \frac{\Gamma(\alpha+\beta)}{\Gamma(\alpha)\Gamma(\beta)}\\
            & = \frac{\alpha}{\alpha+\beta}
            \label{eq:mean-beta-var-alt}
        \end{align}
        onde usamos a defini\'c\~ao da fun\'c\~ao Beta em termos da fun\'c\~ao Gamma e a propriedade $\Gamma(t+1) = t \Gamma(t)$.
    \end{solution}
    
    \item  \label{q:bernoulli-posterior-predictive-alt} Mostre que a probabilidade posterior preditiva $p( x=1 | \data )$ para um
    novo ponto de dado $x$ observado de forma independente \'e igual \`a m\'edia posterior
    de $p(\theta | \data)$, que por sua vez \'e dada por
    \begin{align}
        \E( \theta | \data) & = \frac{\alpha_0 + n_{x=1}}{\alpha_0+\beta_0+n}.
        \label{eq:bernoulli-posterior-mean-alt}
    \end{align}
    \begin{solution}
        Obtemos
        \begin{align}
            p(x=1 | \data) & = \int_0^1 p(x=1, \theta | \data) \ud \theta \quad \quad \text{(regra da soma)}\\
            & = \int_0^1 p(x=1 | \theta, \data) p(\theta | \data) \ud \theta \quad \quad \text{(regra do produto)} \\
            & = \int_0^1 p(x=1 | \theta) p(\theta | \data) \ud \theta \quad \quad (x \independent \data | \theta) \\
            & = \int_0^1 \theta p(\theta | \data) \ud \theta\\
            & = \E[ \theta | \data]
        \end{align}
        Da quest\~ao anterior, conhecemos a m\'edia de uma vari\'avel aleat\'oria Beta. Como $\theta \sim \BetaDist(\theta ; \alpha_n, \beta_n)$, obtemos
        \begin{align}
            p(x=1 | \data) & =  \E[ \theta | \data]\\
            & = \frac{\alpha_n}{\alpha_n+\beta_n}\\
            & = \frac{\alpha_0 +  n_{x=1}} {\alpha_0+ n_{x=1} + \beta_0 +  n_{x=0}}\\
            & = \frac{\alpha_0 +  n_{x=1}}{\alpha_0+ \beta_0 + n}
        \end{align}
        onde a \'ultima equa\'c\~ao decorre do fato de que $n = n_{x=0}+n_{x=1}$. Note que para $n \to \infty$, a m\'edia posterior tende ao MLE
        $n_{x=1}/n$.
        
    \end{solution}
    
    
\end{exenumerate}

% --------------------------------------------------

\ex{Infer\^encia Bayesiana de tabelas de probabilidade em modelos gr\'aficos direcionados totalmente observados de vari\'aveis bin\'arias}
\label{ex:Bayesian-Inference-DGM-alt}

Este \'e o an\'alogo Bayesiano do Exerc\'icio \ref{ex:MLE-DGM-alt} e a
nota\'c\~ao segue aquele exerc\'icio. Consideramos o modelo Bayesiano
\begin{align}
    p(\x | \thetab) &= \prod_{i=1}^d p(x_i | \pa_i, \thetab_i) \quad \quad x_i \in \{0,1\} \label{eq:p-dag2-alt}\\
    p(\thetab; \alphab_0, \betab_0) & = \prod_{i=1}^d \prod_{s=1}^{S_i} \BetaDist(\theta_i^s; \alpha_{i,0}^s, \beta_{i,0}^s)
\end{align}
onde $p(x_i | \pa_i, \thetab_i)$ \'e definido via \eqref{eq:pcond-fac-alt}, $\alphab_0$ \'e um vetor de hiperpar\^ametros contendo todos os
$\alpha_{i,0}^s$, $\betab_0$ o vetor contendo todos os
$\beta_{i,0}^s$, e como antes $\BetaDist$ denota a
distribui\'c\~ao Beta. Sob a priori, todos os par\^ametros s\~ao independentes.

\begin{exenumerate}
    \item Para dados iid $\data = \{\x^{(1)}, \ldots, \x^{(n)}\}$ mostre que 
    \begin{align}
        p(\thetab | \data) & = \prod_{i=1}^d \prod_{s=1}^{S_i} \BetaDist(\theta_i^s, \alpha_{i,n}^s, \beta_{i,n}^s)
    \end{align}
    onde
    \begin{align}
        \alpha_{i,n}^s & = \alpha_{i,0}^s + n_{x_i=1}^s &  \beta_{i,n}^s & = \beta_{i,0}^s + n_{x_i=0}^s
    \end{align}
    e que os par\^ametros tamb\'em s\~ao independentes sob a posterior.
    
    \begin{solution}
        Come\'camos com
        \begin{equation}
            p(\thetab | \data) \propto p(\data | \thetab) p(\thetab; \alphab_0, \betab_0).
        \end{equation}
        Inserir a express\~ao para $p(\data | \thetab)$ dada em \eqref{eq:joint-dag-alt} e a forma assumida da priori resulta em
        \begin{align}
            p(\thetab | \data) &\propto  \prod_{i=1}^d \prod_{s=1}^{S_i}  (\theta_i^s)^{n_{x_i=1}^s} (1-\theta_i^s)^{n_{x_i=0}^s} \prod_{i=1}^d \prod_{s=1}^{S_i} \BetaDist(\theta_i^s; \alpha_{i,0}^s, \beta_{i,0}^s)\\
            & \propto \prod_{i=1}^d \prod_{s=1}^{S_i}   (\theta_i^s)^{n_{x_i=1}^s} (1-\theta_i^s)^{n_{x_i=0}^s}  \BetaDist(\theta_i^s; \alpha_{i,0}^s, \beta_{i,0}^s)\\
            & \propto \prod_{i=1}^d \prod_{s=1}^{S_i}   (\theta_i^s)^{n_{x_i=1}^s} (1-\theta_i^s)^{n_{x_i=0}^s} (\theta_i^s)^{\alpha_{i,0}^s-1} (1-\theta_i^s)^{\beta_{i,0}^s-1}\\
            & \propto \prod_{i=1}^d \prod_{s=1}^{S_i}   (\theta_i^s)^{\alpha_{i,0}^s+n_{x_i=1}^s-1} (1-\theta_i^s)^{\beta_{i,0}^s+n_{x_i=0}^s-1}\\
            & \propto \prod_{i=1}^d \prod_{s=1}^{S_i} \BetaDist(\theta_i^s; \alpha_{i,0}^s+n_{x_i=1}^s, \beta_{i,0}^s+n_{x_i=0}^s)
        \end{align}
        Pode-se verificar imediatamente que $\BetaDist(\theta_i^s; \alpha_{i,0}^s+n_{x_i=1}^s, \beta_{i,0}^s+n_{x_i=0}^s)$ \'e proporcional \`a marginal $p(\theta_i^s | \data)$, de modo que os par\^ametros s\~ao independentes sob a posterior tamb\'em.
    \end{solution}
    
    \item \label{q:dgm-posterior-predictive-alt} Para uma vari\'avel $x_i$ com pais $\pa_i$, compute a probabilidade posterior preditiva $p(x_i=1|\pa_i, \data)$
    
    \begin{solution}
        A solu\'c\~ao \'e an\'aloga \`a solu\'c\~ao da quest\~ao
        \ref{q:bernoulli-posterior-predictive-alt}, usando a regra da soma,
        independ\^encias e propriedades de vari\'aveis aleat\'orias Beta:
        \begin{align}
            p(x_i=1 |\pa_i=s, \data) &= \int p(x_i=1, \theta_i^s | \pa_i=s,
            \data) \ud \theta_i^s\\ & = \int p(x_i=1 | \theta_i^s, \pa_i=s,
            \data) p(\theta_i^s | \pa_i=s , \data)\\ & = \int p(x_i=1 |
            \theta_i^s, \pa_i=s) p(\theta_i^s | \data)\\ & = \int \theta_i^s
            p(\theta_i^s | \data)\\ & = \E[ \theta_i^s | \data)]\\ &
            \overset{\eqref{eq:mean-beta-var-alt}}{=}
            \frac{\alpha_{i,n}^s}{\alpha_{i,n}^s+\beta_{i,n}^s}\\ & =
            \frac{\alpha_{i,0}^s+n^s_{x_i=1}}{\alpha_{i,0}^s+\beta_{i,0}^s+n^s}
            \label{eq:dgm-posterior-predictive-alt}
        \end{align}
    \end{solution}
    onde $n^s = n^s_{x_i=0}+n^s_{x_i=1}$ denota o n\'umero de vezes que
    a configura\'c\~ao dos pais $s$ ocorre nos dados observados $\data$.
    
\end{exenumerate}

\ex{Exemplo c\^ancer-amianto-fumo: infer\^encia Bayesiana}

Considere o modelo especificado pelo DAG

\begin{center}
    \begin{tikzpicture}[dgraph]
        \node[cont] (x) at (0,0) {a};
        \node[cont] (y) at (2,0) {s};
        \node[cont] (z) at (1,-1) {c};
        \draw(x) -- (z);
        \draw(y) -- (z);
    \end{tikzpicture}
\end{center}

As distribui\'c\~oes de $a$ e $s$ s\~ao distribui\'c\~oes de Bernoulli com
par\^ametro (probabilidade de sucesso) $\theta_a$ e $\theta_s$,
respectivamente, i.e.\
\begin{equation}
    p(a | \theta_a) = \theta_a^a(1-\theta_a)^{1-a} \quad \quad p(s | \theta_s) = \theta_s^s(1-\theta_s)^{1-s},
\end{equation}
e a distribui\'c\~ao de $c$ dados os pais \'e parametrizada conforme especificado na tabela seguinte
\begin{center}
    \begin{tabular}{@{}rll@{}}
        \toprule
        $p(c=1 | a,s, \theta^1_c, \ldots, \theta_c^4))$ & $a$ & $s$\\
        \midrule
        $\theta^1_c$ & 0 & 0\\
        $\theta^2_c$ & 1 & 0\\
        $\theta^3_c$ & 0 & 1\\
        $\theta^4_c$ & 1 & 1\\
        \bottomrule
    \end{tabular}
\end{center}
Assumimos que a priori sobre os par\^ametros do modelo, $(\theta_a,
\theta_s, \theta^1_c, \ldots, \theta^4_c)$, se fatoriza e \'e dada por
distribui\'c\~oes Beta com hiperpar\^ametros $\alpha_0 = 1$ e $\beta_0 = 1$
(iguais para todos os par\^ametros).

Assuma que observamos os seguintes dados iid (cada linha \'e um ponto de dado).

\begin{center}
    \scalebox{1}{
        \begin{tabular}{lll}
            \toprule
            a & s & c\\
            \midrule
            0 &   1 &   1\\
            0 &   0 &   0\\
            1 &   0 &   1\\
            0 &   0 &   0\\
            0 &   1 &   0\\
            \bottomrule
    \end{tabular}}
\end{center}

\begin{exenumerate}
    \item Determine as probabilidades preditivas posteriores $p(a =1 |\data)$ e $p(s=1|\data)$. 
    \begin{solution}
        Com o \exref {ex:Bayesian-inference-Bernoulli-alt} quest\~ao
        \ref{q:bernoulli-posterior-predictive-alt}, temos
        \begin{align}
            p(a =1 |\data) &= \E(\theta^a | \data) = \frac{1 + 1}{1+1+5} = \frac{2}{7}\\
            p(s=1|\data) &=\E(\theta^s | \data) = \frac{1 + 2}{1+1+5} = \frac{3}{7}
        \end{align}
    \end{solution}
    
    \item Determine as probabilidades preditivas posteriores $p(c=1 | \pa,
    \data)$ para todas as configura\'c\~oes de pais poss\'iveis.
    
    \begin{solution}
        Os pais de $c$ s\~ao $(a,s)$. Com o
        \exref{ex:Bayesian-Inference-DGM-alt} quest\~ao
        \ref{q:dgm-posterior-predictive-alt}, temos
        \begin{center}
            \begin{tabular}{@{}lll@{}}
                \toprule
                $p(c=1 | a,s, \data)$ & $a$ & $s$\\
                \midrule
                $(1+0)/(1+1+2)=1/4 $ & 0 & 0\\
                $(1+1)/(1+1+1)=2/3 $ & 1 & 0\\
                $(1+1)/(1+1+2)=1/2 $ & 0 & 1\\
                $(1+0)/(1+1)=1/2 $ & 1 & 1\\
                \bottomrule
            \end{tabular}
        \end{center}
        Comparado \`a solu\'c\~ao MLE na
        \exref{q:cancer-smoking-asbestos-mle-alt-q} quest\~ao
        \ref{q:cancer-smoking-asbestos-mle-alt-q}, vemos que as estimativas s\~ao
        menos extremas. Isso ocorre porque s\~ao uma combina\'c\~ao do
        conhecimento a priori e dos dados observados. Al\'em disso, quando n\~ao temos
        nenhum dado, a posterior \'e igual \`a priori, ao contr\'ario do MLE, onde a
        estimativa n\~ao est\'a definida.
        
    \end{solution}
    
\end{exenumerate}

% --------------------------------------------------


\ex{Aprendizado de par\^ametros de um modelo gr\'afico direcionado}

Consideramos o modelo gr\'afico direcionado mostrado abaixo \`a esquerda para
as quatro vari\'aveis bin\'arias $t,b,s,x$, cada uma sendo zero ou
um. Assuma que observamos os dados mostrados na tabela \`a
direita.\\

\begin{minipage}[t]{0.45 \textwidth}
    \begin{center}
        {\small Modelo:\\[2ex]}
        \scalebox{1}{ % x,y
            \begin{tikzpicture}[dgraph]
                \node[cont] (t) at (0,0) {$t$};
                \node[cont] (b) at (2,0) {$b$};
                \node[cont] (s) at (1,-1.5) {$s$};
                \node[cont] (x) at (-1,-1.5) {$x$};
                
                \draw (t) -- (s);
                \draw (b) -- (s);
                \draw (t) -- (x);
        \end{tikzpicture}}
    \end{center}\vspace{2ex}
    \begin{tabular}{l l}
        $t=1$& tem tuberculose\\
        $b=1$& tem bronquite \\
        $s=1$& tem falta de ar\\
        $x=1$& tem raio-x positivo
    \end{tabular}
\end{minipage}
\hspace{2ex}
\begin{minipage}[t]{0.45\textwidth}
    \begin{center}
        {\small Dados observados:\\[2ex]}
        \scalebox{1}{
            \begin{tabular}{llll}
                \toprule
                x & s & t & b\\
                \midrule
                0 &   1 &   0 &   1\\
                0 &   0 &   0 &   0\\
                0 &   1 &   0 &   1\\
                0 &   1 &   0 &   1\\
                0 &   0 &   0 &   0\\
                0 &   0 &   0 &   0\\
                0 &   1 &   0 &   1\\
                0 &   1 &   0 &   1\\
                0 &   0 &   0 &   1\\
                1 &   1 &   1 &   0\\
                \bottomrule
        \end{tabular}}
    \end{center}
\end{minipage}\\[2ex]
Assumimos que a fmp (condicional) de $s|t,b$ \'e especificada pela seguinte
tabela de probabilidade parametrizada:
\begin{center}
    \begin{tabular}{@{}cll@{}}
        \toprule
        $p(s=1 | t, b; \theta^1_s, \ldots, \theta_s^4))$ & $t$ & $b$\\
        \midrule
        $\theta^1_s$ & 0 & 0\\
        $\theta^2_s$ & 1 & 0\\
        $\theta^3_s$ & 0 & 1\\
        $\theta^4_s$ & 1 & 1\\
        \bottomrule
    \end{tabular}
\end{center}

\begin{exenumerate}
    \item Quais s\~ao as estimativas de m\'axima verossimilhan\'ca para $p(s=1 | b=0,
    t=0)$ e $p(s=1 | b=0, t=1)$, i.e.\ os par\^ametros $\theta^1_s$
    e $\theta^2_s$?
    
    \begin{solution}
        As estimativas de m\'axima verossimilhan\'ca (MLEs) s\~ao iguais \`a fra\'c\~ao de ocorr\^encias dos eventos relevantes.
        \begin{align}
            \hat{\theta}^1_s & =\frac{\sum_{i=1}^n \ind(s_i=1, b_i=0, t_i=0)}{\sum_{i=1}^n \ind(b_i=0, t_i=0)} = \frac{0}{3} = 0\\
            \hat{\theta}^2_s & =\frac{\sum_{i=1}^n \ind(s_i=1, b_i=0, t_i=1)}{\sum_{i=1}^n \ind(b_i=0, t_i=1)} = \frac{1}{1} = 1
        \end{align}
        
    \end{solution}
    
    \item Assuma que cada par\^ametro na tabela para $p(s | t,b)$ possua uma
    priori uniforme em $(0,1)$. Compute a m\'edia posterior dos
    par\^ametros de \mbox{$p(s=1 | b=0, t=0)$} e $p(s=1 | b=0, t=1)$
    e explique a diferen\'ca em rela\'c\~ao \`as estimativas de m\'axima verossimilhan\'ca.
    
    \begin{solution}
        Uma priori uniforme corresponde a uma distribui\'c\~ao Beta com
        hiperpar\^ametros $\alpha_0=\beta_0=1$. Com o
        \exref{ex:Bayesian-Inference-DGM-alt} quest\~ao
        \ref{q:dgm-posterior-predictive-alt}, temos
        \begin{align}
            \E(\theta_s^1 | \data) & = \frac{\alpha_0+0}{\alpha_0+\beta_0+3} = \frac{1}{5}\\
            \E(\theta_s^2 | \data) & = \frac{\alpha_0+1}{\alpha_0+\beta_0+1} = \frac{2}{3}
        \end{align}
        Comparada ao MLE, a m\'edia posterior \'e menos extrema. Pode ser
        considerada uma estimativa ``suavizada'' ou regularizada, onde
        $\alpha_0 > 0$ e $\beta_0 > 0$ fornecem a regulariza\'c\~ao (veja
        \url{https://en.wikipedia.org/wiki/Additive_smoothing}). Podemos
        perceber uma atra\'c\~ao dos par\^ametros em dire\'c\~ao \`a m\'edia preditiva a priori,
        que \'e igual a 1/2.
        
    \end{solution}
    
\end{exenumerate}

% --------------------------------------------------------------- 

\ex{An\'alise de fatores}
\label{ex:FA-alt}

Um amigo prop\~oe melhorar o modelo de an\'alise de fatores
trabalhando com vari\'aveis latentes correlacionadas. O modelo proposto \'e
\begin{align}
    p(\h; \C) &= \Gauss(\h; \zerob, \C) & p(\v | \h; \F, \Psib, \c) = \Gauss(\v; \F \h+\c, \Psib)
\end{align}
onde $\C$ \'e uma matriz de covari\^ancia $H \times H$, $\F$ \'e a matriz $D
\times H$ com as cargas fatoriais (*factor loadings*), $\Psib=\diag(\Psi_1,
\ldots, \Psi_D)$, $\c\in \mathbb{R}^D$ e a dimens\~ao das
latentes $H$ \'e menor que a dimens\~ao das vis\'iveis
$D$. $\Gauss(\x; \mub, \Sigmab)$ denota a fdp de uma Gaussiana com
m\'edia $\mub$ e matriz de covari\^ancia $\Sigmab$. O modelo de an\'alise de fatores padr\~ao \'e obtido quando $\C$ \'e a matriz identidade.

\begin{exenumerate}
    
    \item Qual \'e a distribui\'c\~ao marginal das vis\'iveis $p(\v; \thetab)$, onde $\thetab$ representa os par\^ametros $\C, \F, \c, \Psib$?
    
    \begin{solution}
        As especifica\'c\~oes do modelo s\~ao equivalentes ao seguinte processo de gera\'c\~ao de dados:
        \begin{align}
            \h &\sim \Gauss(\h; \zerob, \C) & \epsilonb &\sim \Gauss(\epsilonb; \zerob, \Psib) &  \v &= \F \h + \c + \epsilonb
        \end{align}
        Recorde o resultado b\'asico sobre a distribui\'c\~ao de transforma\'c\~oes lineares
        de Gaussianas: se $\x$ possui densidade $\Gauss(\x;
        \mub_x, \C_x)$, $\z$ densidade $\Gauss(\z; \mub_z, \C_z)$, e $\x
        \independent \z$ ent\~ao $\y = \A \x + \z$ possui densidade
        $$\Gauss(\y; \A\mub_x+\mub_z, \A \C_x \A^\top + \C_z).$$
        Segue, portanto, que $\v$ \'e Gaussiana com m\'edia $\mub$ e covari\^ancia $\Sigmab$,
        \begin{align}
            \mub &= \F \underbrace{\E[\h]}_{\zerob} + \c + \underbrace{\E[\epsilonb]}_{\zerob}\\
            & = \c\\
            \Sigmab &= \F \Var[\h] \F^\top + \Var[\epsilonb]\\
            & = \F \C \F^\top + \Psib.
        \end{align}
        
        
    \end{solution}
    
    \item Assuma que a decomposi\'c\~ao em valores singulares de $\C$ seja dada por
    \begin{align}
        \C & = \EE \Lambdab \EE^\top
    \end{align}
    onde $\Lambdab = \diag(\lambda_1, \ldots, \lambda_D)$ \'e uma matriz diagonal
    contendo os autovalores, e $\EE$ \'e uma matriz ortonormal
    contendo os autovetores correspondentes. A raiz quadrada matricial de $\C$ \'e a matriz $\M$ tal que
    \begin{align}
        \M \M = \C,
    \end{align}
    e n\'os a denotamos por $\C^{1/2}$. Mostre que a raiz quadrada matricial de
    $\C$ \'e igual a
    \begin{align}
        \C^{1/2} = \EE \diag(\sqrt{\lambda_1}, \ldots, \sqrt{\lambda_D}) \EE^\top.
    \end{align}
    
    \begin{solution}
        Verificamos que $\C^{1/2}\C^{1/2} = \C$:
        \begin{align}
            \C^{1/2} \C^{1/2} & = \EE \diag(\sqrt{\lambda_1}, \ldots, \sqrt{\lambda_D})\EE^\top \EE \diag(\sqrt{\lambda_1}, \ldots, \sqrt{\lambda_D})\EE^\top\\
            & = \EE \diag(\sqrt{\lambda_1}, \ldots, \sqrt{\lambda_D}) \; \I \; \diag(\sqrt{\lambda_1}, \ldots, \sqrt{\lambda_D})\EE^\top\\
            &= \EE \diag(\sqrt{\lambda_1}, \ldots, \sqrt{\lambda_D}) \diag(\sqrt{\lambda_1}, \ldots, \sqrt{\lambda_D})\EE^\top\\
            &= \EE \diag(\lambda_1, \ldots, \lambda_D) \EE^\top\\
            &= \EE \Lambdab \EE^\top\\
            & =\C
        \end{align}
        
    \end{solution}
    
    \item Mostre que o modelo de an\'alise de fatores proposto \'e equivalente ao modelo de an\'alise de fatores original
    \begin{align}
        p(\h; \I) &= \Gauss(\h; \zerob, \I) & p(\v | \h; \tilde{\F}, \Psib, \c) = \Gauss(\v; \tilde{\F} \h+\c, \Psib)
    \end{align}
    com $\tilde{\F} = \F \C^{1/2}$, de modo que os par\^ametros extras fornecidos
    pela matriz de covari\^ancia $\C$ s\~ao, na verdade, redundantes e nada \'e ganho com a parametriza\'c\~ao mais rica. 
    
    \begin{solution}
        Verificamos que o modelo possui a mesma distribui\'c\~ao para as vis\'iveis. Como antes $\E[\v] = \c$, e a matriz de covari\^ancia \'e
        \begin{align}
            \Var[\v] & = \tilde{\F} \I \tilde{\F}^\top + \Psib\\
            & =  \F \C^{1/2} \C^{1/2}  \F^\top + \Psib\\
            & =  \F \C  \F^\top + \Psib
        \end{align}
        onde usamos que $\C^{1/2}$ \'e uma matriz sim\'etrica. Isso significa
        que a correla\'c\~ao entre os $\h$ pode ser absorvida pela
        matriz de fatores $\F$ e o conjunto de fdps definidas pelo modelo
        proposto \'e igual ao conjunto de fdps do modelo de an\'alise de fatores original.
        
        Outra forma de ver o resultado \'e considerar o processo de gera\'c\~ao de
        dados e notar que podemos amostrar $\h$ de $\Gauss(\h; \zerob, \C)$
        primeiro amostrando $\h'$ de $\Gauss(\h'; \zerob, \I)$ e depois transformando a amostra por $\C^{1/2}$,
        \begin{align}
            \h &\sim \Gauss(\h; \zerob, \C)  & \Longleftrightarrow && \h &= \C^{1/2} \h' \quad \quad \quad \h' \sim \Gauss(\h'; \zerob, \I).
        \end{align}
        Isso decorre novamente das propriedades b\'asicas das transforma\'c\~oes lineares de Gaussianas, i.e.
        $$ \Var(\C^{1/2} \h') = \C^{1/2}\Var(\h')(\C^{1/2})^\top =  \C^{1/2}\I \C^{1/2} =  \C$$
        e $\E(\C^{1/2} \h') = \C^{1/2} \E(\h') = \zerob$.
        
        Para gerar amostras do modelo de an\'alise de fatores proposto, proceder\'iamos da seguinte forma:
        \begin{align}
            \h' &\sim \Gauss(\h'; \zerob, \I) & \epsilonb &\sim \Gauss(\epsilonb; \zerob, \Psib) &  \v &= \F (\C^{1/2}\h') + \c + \epsilonb
        \end{align}
        Mas o termo
        $$ \v = \F (\C^{1/2}\h') + \c + \epsilonb$$
        pode ser escrito como
        $$\v = (\F \C^{1/2} ) \h' + \c + \epsilonb = \tilde{\F} \h' + \c + \epsilonb$$
        e como $\h'$ segue $\Gauss(\h'; \zerob, \I)$, estamos de volta ao modelo de an\'alise de fatores original.
        
    \end{solution}
\end{exenumerate}



\ex{An\'alise de componentes independentes (ICA)}
\label{ex:ex2-alt}

\begin{exenumerate}
    
    \item O branqueamento (*whitening*) corresponde a transformar linearmente uma vari\'avel aleat\'oria
    $\x$ (ou os dados correspondentes) de modo que a vari\'avel aleat\'oria resultante $\z$ possua uma matriz de covari\^ancia identidade, i.e.\ 
    $$\z = \V \x \quad \text{com} \quad \Var[\x] = \C \quad
    \text{e}\quad \Var[\z] = \I. $$ A matriz $\V$ \'e chamada de
    matriz de branqueamento. N\~ao fazemos uma hip\'otese distribucional sobre
    $\x$; em particular, $\x$ pode ou n\~ao ser Gaussiana.
    
    Dada a decomposi\'c\~ao em autovalores $\C = \EE \Lambdab \EE^\top$,
    mostre que
    \begin{equation}
        \V = \diag (\lambda_1^{-1/2}, \ldots,  \lambda_d^{-1/2}) \EE^\top
    \end{equation}
    \'e uma matriz de branqueamento.
    
    \begin{solution}
        De $\Var[\z] = \Var[\V\x]= \V \Var[\x] \V^\top$, segue que
        \begin{align}
            \Var[\z] & =  \V \Var[\x] \V^\top\\
            & =  \V \C \V^\top\\
            & =  \V \EE \Lambdab \EE^\top \V^\top\\
            & =   \diag (\lambda_1^{-1/2}, \ldots,
            \lambda_d^{-1/2}) \EE^\top \EE \Lambdab \EE^\top  \V^\top\\
            & = \diag (\lambda_1^{-1/2}, \ldots,
            \lambda_d^{-1/2}) \Lambdab \EE^\top \V^\top
        \end{align}
        onde usamos que $\EE^\top \EE = \I$. Como
        $$ \V^\top = \left[  \diag (\lambda_1^{-1/2}, \ldots, \lambda_d^{-1/2}) \EE^\top \right]^\top = \EE  \diag (\lambda_1^{-1/2}, \ldots, \lambda_d^{-1/2})$$
        temos ainda
        \begin{align}
            \Var[\z] & = \diag (\lambda_1^{-1/2}, \ldots,
            \lambda_d^{-1/2}) \Lambdab \EE^\top  \EE  \diag (\lambda_1^{-1/2}, \ldots, \lambda_d^{-1/2})\\
            & =  \diag (\lambda_1^{-1/2}, \ldots,
            \lambda_d^{-1/2}) \Lambdab \diag (\lambda_1^{-1/2}, \ldots, \lambda_d^{-1/2})\\
            & =  \diag (\lambda_1^{-1/2}, \ldots,
            \lambda_d^{-1/2}) \diag(\lambda_1, \ldots, \lambda_d) \diag (\lambda_1^{-1/2}, \ldots, \lambda_d^{-1/2})\\
            & = \I,
        \end{align}
        o que mostra que $\V$ \'e de fato uma matriz de branqueamento v\'alida. Note que as matrizes de branqueamento n\~ao s\~ao \'unicas. Por exemplo,
        $$ \tilde{\V} = \EE \diag (\lambda_1^{-1/2}, \ldots,
        \lambda_d^{-1/2}) \EE^\top$$ tamb\'em \'e uma matriz de branqueamento
        v\'alida. Mais geralmente, se $\V$ for uma matriz de branqueamento, ent\~ao $\R \V$ tamb\'em ser\'a
        uma matriz de branqueamento quando $\R$ for uma matriz
        ortonormal. Isso ocorre porque
        $$ \Var[ \R \V \x ] = \R \Var[\V \x] \R^\top = \R \I \R^{\top} = \I $$
        onde usamos que $\V$ \'e uma matriz de branqueamento tal que $\V \x$ possui matriz de covari\^ancia identidade.
        
        
    \end{solution}
    
    \item Considere o modelo ICA
    \begin{align}
        \v &= \A \h, & \h &\sim p_{\h}(\h), &  p_{\h}(\h)&= \prod_{i=1}^D p_h(h_i),
    \end{align}
    onde a matriz $\A$ \'e invers\'ivel e os $h_i$ s\~ao vari\'aveis
    aleat\'orias independentes de m\'edia zero e vari\^ancia um. Seja $\V$ uma
    matriz de branqueamento para $\v$. Mostre que $\z = \V \v$ segue o modelo ICA
    \begin{align}
        \z &= \tilde{\A} \h, & \h &\sim p_{\h}(\h), &  p_{\h}(\h)&= \prod_{i=1}^D p_h(h_i),
    \end{align}
    onde $\tilde{\A}$ \'e uma matriz ortonormal.
    
    \begin{solution}
        Se $\v$ segue o modelo ICA, temos
        \begin{align}
            \z & = \V \v\\
            & = \V \A \h\\
            &= \tilde{\A} \h
        \end{align}
        com $\tilde{\A} = \V \A$. Pela opera\'c\~ao de branqueamento, a matriz de covari\^ancia de $\z$ \'e a identidade, de modo que
        \begin{align}
            \I = \Var(\z) =  \tilde{\A} \Var(\h)  \tilde{\A}^{\top}.
        \end{align}
        Pelo modelo ICA, $\Var(\h) = \I$, de modo que $\tilde{\A}$ deve satisfazer
        \begin{equation}
            \I =  \tilde{\A} \tilde{\A}^{\top},
        \end{equation}
        o que significa que $\tilde{\A}$ \'e ortonormal.
        
        No modelo ICA original, o n\'umero de par\^ametros \'e dado pelo
        n\'umero de elementos da matriz $\A$, que \'e $D^2$ se $\v$ for
        D-dimensional. Uma matriz ortogonal possui $D(D-1)/2$ graus de
        liberdade {\small (veja
            e.g. \url{https://en.wikipedia.org/wiki/Orthogonal_matrix}}), de
        modo que podemos pensar que o branqueamento ``resolve metade do problema de ICA''. Como o branqueamento \'e uma opera\'c\~ao padr\~ao relativamente simples, muitos algoritmos \citep[e.g.\ ``fastICA'',][]{Hyvarinen1999}
        primeiro reduzem a complexidade do problema de estimativa branqueando os
        dados. Al\'em disso, devido \`as propriedades da matriz ortogonal,
        a log-verossimilhan\'ca para o modelo ICA tamb\'em se simplifica para
        dados branqueados: a log-verossimilhan\'ca para o modelo ICA sem branqueamento \'e
        \begin{equation}
            \ell(\B) = \sum_{i=1}^n  \sum_{j=1}^D\log p_h (\b_j \v_i) + n \log |\det \B|
        \end{equation}
        onde $\B = \A^{-1}$. Se primeiro branquearmos os dados, a log-verossimilhan\'ca torna-se
        \begin{equation}
            \ell(\tilde{\B}) = \sum_{i=1}^n  \sum_{j=1}^D\log p_h (\tilde{\b}_j \z_i) + n \log |\det \tilde{\B}|
        \end{equation}
        onde $\tilde{\B} = \tilde{\A}^{-1} = \tilde{\A}^\top$ j\'a que $\tilde{\A}$ \'e
        uma matriz ortogonal. Isso significa que $\tilde{\B}^{-1} = \tilde{\A} =
        \tilde{\B}^\top$ e $\tilde{\B}$ \'e uma matriz ortogonal. Assim,
        $\det \tilde{\B}=1$, e o termo $\log \det$ \'e zero. Logo, a
        log-verossimilhan\'ca em dados branqueados simplifica-se para
        \begin{equation}
            \ell(\tilde{\B}) = \sum_{i=1}^n  \sum_{j=1}^D\log p_h (\tilde{\b}_j \z_i).
        \end{equation}
        Embora a log-verossimilhan\'ca assuma uma forma mais simples, o problema de otimiza\'c\~ao \'e agora um problema de otimiza\'c\~ao restrito: $\tilde{\B}$ \'e
        restrita a ser ortonormal. Para mais informa\'c\~oes, veja
        e.g.\ \citep[Cap\'itulo 9]{Hyvarinen2001}.  
        
    \end{solution}
    
\end{exenumerate}


% --------------------------------------------------------------- 

\ex{Score matching para a fam\'ilia exponencial}
\label{ex:score-matching-exp-family-alt}
A fun\'c\~ao objetivo $J(\thetab)$ que \'e minimizada no *score matching* \'e
\begin{align} 
    J(\thetab) &= \frac{1}{n} \sum_{i=1}^n \sum_{j=1}^m \left[ \partial_j
    \psi_j(\x_i;\thetab) + \frac{1}{2} \psi_j(\x_i;\thetab)^2 \right],\label{eq:Jsm-def-alt}
\end{align}
onde $\psi_j$ \'e a derivada parcial da log-fdp do modelo $\log
p(\x ;\thetab)$ em rela\'c\~ao \`a $j$-\'esima coordenada (inclina\'c\~ao) e
$\partial_j \psi_j$ sua segunda derivada parcial (curvatura). Os
dados observados s\~ao denotados por $\x_1,\ldots,\x_n$ e $\x \in
\mathbb{R}^m$.

O objetivo deste exerc\'icio \'e mostrar que, para modelos estat\'isticos da forma
\begin{align}
    \log p(\x;\thetab) &= \sum_{k=1}^K \theta_k F_k(\x) -\log Z(\thetab), \quad \quad \quad \x \in \mathbb{R}^m,
    \label{eq:exp-family-def-alt}
\end{align}
a fun\'c\~ao objetivo do score matching torna-se uma forma quadr\'atica, que pode ser otimizada eficientemente \citep[veja e.g.\ ][Ap\^endice A.5.3]{Barber2012}.

O conjunto de modelos acima \'e chamado de fam\'ilia exponencial (cont\'inua), ou tamb\'em de modelos log-lineares, porque os modelos s\~ao lineares nos
par\^ametros $\theta_k$. Como a fam\'ilia exponencial geralmente inclui
fun\'c\~oes de massa de probabilidade tamb\'em, o qualificador ``cont\'inua'' pode
ser usado para destacar que estamos considerando aqui apenas vari\'aveis aleat\'orias cont\'inuas. As fun\'c\~oes $F_k(\x)$ s\~ao assumidas conhecidas (elas
s\~ao chamadas de estat\'isticas suficientes).


\begin{exenumerate}
    
    \item Denomine por $\K(\x)$ a matriz com elementos $K_{kj}(\x)$,
    \begin{equation}
        K_{kj}(\x) = \frac{\partial F_k(\x)}{\partial x_j}, \quad \quad \quad k=1 \ldots K, \quad j=1 \ldots m,
    \end{equation}
    e por $\H(\x)$ a matriz com elementos $H_{kj}(\x)$,
    \begin{equation}
        H_{kj}(\x) = \frac{\partial^2 F_k(\x)}{\partial x_j^2}, \quad \quad \quad k=1 \ldots K, \quad j=1 \ldots m.
    \end{equation}
    Al\'em disso, seja $\h_j(\x) = (H_{1j}(\x), \ldots, H_{Kj}(\x))^\top$ o $j$-\'esimo vetor coluna de $\H(\x)$.
    
    Mostre que para a fam\'ilia exponencial cont\'inua, o objetivo do score matching na Equa\'c\~ao \eqref{eq:Jsm-def-alt} torna-se
    \begin{equation}
        J(\thetab) =  \thetab^\top \r +  \frac{1}{2} \thetab^\top \M \thetab,
    \end{equation}
    onde
    \begin{align}
        \label{eq:Jsm-sol-alt}
        \r&=\frac{1}{n}\sum_{i=1}^n \sum_{j=1}^m \h_j(\x_i), & \M &= \frac{1}{n}\sum_{i=1}^n \K(\x_i)\K(\x_i)^\top.
    \end{align}
    
    \begin{solution}
        Para
        \begin{align}
            \log p(\x;\thetab) &= \sum_{k=1}^K \theta_k F_k(\x) -\log Z(\thetab)
        \end{align}
        a primeira derivada em rela\'c\~ao a $x_j$, o $j$-\'esimo elemento de $\x$, \'e
        \begin{align}
            \psi_j(\x; \thetab) &= \frac{\partial \log p(\x; \thetab)}{\partial x_j}\\
            & = \sum_{k=1}^K \theta_k \frac{\partial F_k(\x)}{\partial x_j}\\
            & =  \sum_{k=1}^K \theta_k K_{kj}(\x).
        \end{align}
        A segunda derivada \'e
        \begin{align}
            \partial_j \psi_j(\x; \thetab) & = \frac{\partial^2 \log p(\x; \thetab)}{\partial x_j^2}\\
            & =  \sum_{k=1}^K \theta_k \frac{\partial^2 F_k(\x)}{\partial x_j^2}\\
            & =  \sum_{k=1}^K \theta_k  H_{kj}(\x),
        \end{align}
        que podemos escrever de forma mais compacta como
        \begin{align}
            \partial_j  \psi_j(\x; \thetab) &= \thetab^\top \h_j(\x).
        \end{align}
        
        O objetivo do score matching na Equa\'c\~ao \eqref{eq:Jsm-def-alt} apresenta a soma $\sum_j \psi_j(\x;\thetab)^2$. O termo $\psi_j(\x;\thetab)^2$ \'e igual a
        \begin{align}
            \psi_j(\x; \thetab)^2 & = \left[ \sum_{k=1}^K \theta_k K_{kj}(\x)\right]^2\\
            & =  \sum_{k=1}^K \sum_{k'=1}^K  K_{kj}(\x) K_{k'j}(\x) \theta_{k} \theta_{k'},
        \end{align}
        de modo que
        \begin{align}
            \sum_{j=1}^m \psi_j(\x; \thetab)^2  &= \sum_{j=1}^m  \sum_{k=1}^K \sum_{k'=1}^K  K_{kj}(\x) K_{k'j}(\x) \theta_{k} \theta_{k'}\\
            & = \sum_{k=1}^K \sum_{k'=1}^K \theta_{k} \theta_{k'} \left[ \sum_{j=1}^m K_{kj}(\x) K_{k'j}(\x) \right],
        \end{align}
        que pode ser expresso de forma mais compacta usando nota\'c\~ao matricial. Notando
        que $$\sum_{j=1}^m K_{kj}(\x_i) K_{k'j}(\x_i)$$ \'e igual ao
        elemento $(k,k')$ do produto matriz-matriz $\K(\x_i)\K(\x_i)^\top$,
        \begin{equation}
            \sum_{j=1}^m K_{kj}(\x_i) K_{k'j}(\x_i) = \left[\K(\x_i)\K(\x_i)^\top\right]_{k,k'},
        \end{equation}
        podemos escrever
        \begin{align}
            \sum_{j=1}^m \psi_j(\x; \thetab)^2 & =  \sum_{k=1}^K \sum_{k'=1}^K \theta_{k} \theta_{k'}  \left[\K(\x_i)\K(\x_i)^\top\right]_{k,k'}\\
            & =  \thetab^\top \K(\x_i)\K(\x_i)^\top \thetab
        \end{align}
        onde usamos que para alguma matriz $\A$
        \begin{equation}
            \thetab^\top \A \thetab = \sum_{k,k'} \theta_k \theta_{k'} [\A]_{k,k'}
        \end{equation}
        onde $[\A]_{k,k'}$ \'e o elemento $(k,k')$ da matriz $\A$.
        
        Inserir as express\~oes na Equa\'c\~ao \eqref{eq:Jsm-def-alt} resulta em
        \begin{align} 
            J(\thetab) &= \frac{1}{n} \sum_{i=1}^n  \sum_{j=1}^m\left[ \partial_j \psi_j(\x_i;\thetab) + \frac{1}{2} \psi_j(\x_i;\thetab)^2 \right] \\
            &= \frac{1}{n} \sum_{i=1}^n  \sum_{j=1}^m \partial_j \psi_j(\x_i;\thetab) + \frac{1}{2}  \frac{1}{n} \sum_{i=1}^n  \sum_{j=1}^m \psi_j(\x_i;\thetab)^2 \\
            & =  \frac{1}{n} \sum_{i=1}^n \sum_{j=1}^m \thetab^\top \h_j(\x_i) + \frac{1}{2} \frac{1}{n} \sum_{i=1}^n \thetab^\top \K(\x_i)\K(\x_i)^\top \thetab\\
            & =  \thetab^\top \left[  \frac{1}{n} \sum_{i=1}^n \sum_{j=1}^m \h_j(\x_i)\right] + \frac{1}{2}  \thetab^\top \left[ \frac{1}{n}\sum_{i=1}^n  \K(\x_i)\K(\x_i)^\top \right] \thetab\\
            & = \thetab^\top \r + \frac{1}{2}\thetab^\top \M \thetab,
        \end{align}
        que \'e o resultado desejado. 
    \end{solution}
    
    
    \item A fdp de uma Gaussiana de m\'edia zero parametrizada pela vari\^ancia $\sigma^2$ \'e 
    \begin{align}
        p(x; \sigma^2) = \frac{1}{\sqrt{2 \pi \sigma^2}}\exp\left(-\frac{x^2}{2\sigma^2} \right), \quad \quad \quad x \in \mathbb{R}.
    \end{align}
    A Gaussiana (multivariada) \'e um membro da fam\'ilia exponencial.
    Por compara\'c\~ao com a Equa\'c\~ao \eqref{eq:exp-family-def-alt}, podemos
    reparametrizar o modelo estat\'istico $\{p(x; \sigma^2)\}_{\sigma^2}$ e trabalhar com
    \begin{align}
        p(x; \theta) = \frac{1}{Z(\theta)}\exp\left(\theta x^2 \right), \quad \quad \theta < 0, \quad \quad x \in \mathbb{R},
    \end{align}
    em vez disso. As duas parametriza\'c\~oes est\~ao relacionadas por $\theta =
    -1/(2 \sigma^2)$. Usando o resultado anterior sobre a fam\'ilia
    exponencial (cont\'inua), determine a estimativa de score matching
    $\hat{\theta}$, e mostre que a correspondente $\hat{\sigma}^2$ \'e
    a mesma que a estimativa de m\'axima verossimilhan\'ca. Este resultado \'e
    not\'avel porque, ao contr\'ario da estimativa de m\'axima verossimilhan\'ca, o
    score matching n\~ao necessita da fun\'c\~ao de parti\'c\~ao $Z(\theta)$ para a estimativa.
    
    \begin{solution}
        Por compara\'c\~ao com a Equa\'c\~ao \eqref{eq:exp-family-def-alt}, a
        estat\'istica suficiente $F(x)$ \'e $x^2$.
        
        Primeiro determinamos a fun\'c\~ao objetivo do score matching. Para
        isso, precisamos determinar as quantidades $\r$ e $\M$ na
        Equa\'c\~ao \eqref{eq:Jsm-sol-alt}. Aqui, tanto $\r$ quanto $\M$ s\~ao escalares,
        e o mesmo vale para as matrizes $\K$ e $\H$ que definem $\r$ e $\M$. Por
        suas defini\'c\~oes, obtemos
        \begin{align}
            K(x) &=  \frac{\partial F(x)}{\partial x} = 2x\\
            H(x) &= \frac{\partial^2 F(x)}{\partial x^2} = 2\\
            r & = 2\\
            M &=  \frac{1}{n}\sum_{i=1}^n K(x_i)^2 \\
            & = 4 m_2
        \end{align}
        onde $m_2$ denota o segundo momento emp\'irico,
        \begin{equation}
            m_2 = \frac{1}{n} \sum_{i=1}^n x_i^2.
        \end{equation}
        Com a Equa\'c\~ao \eqref{eq:Jsm-def-alt}, o objetivo do score matching \'e portanto
        \begin{align}
            J(\theta) &= 2 \theta + \frac{1}{2} 4 m_2 \theta^2 \\
            & = 2 \theta + 2 m_2 \theta^2
        \end{align}
        Uma condi\'c\~ao necess\'aria para o minimizador satisfazer \'e
        \begin{align}
            \frac{\partial J(\theta)}{\partial \theta} & = 2 + 4 \theta m_2 \\
            & = 0
        \end{align}
        O \'unico valor de par\^ametro que satisfaz a condi\'c\~ao \'e
        \begin{equation}
            \hat{\theta} = -\frac{1}{2 m_2}.
        \end{equation}
        A segunda derivada de $J(\theta)$ \'e
        \begin{align}
            \frac{\partial^2 J(\theta)}{\partial \theta^2} & = 4 m_2,
        \end{align}
        que \'e positiva (desde que nem todos os pontos de dados sejam zero). Assim, $\hat{\theta}$ \'e um minimizador.
        
        Da rela\'c\~ao $\theta = -1/(2 \sigma^2)$, obtemos que a estimativa de score matching da vari\^ancia $\sigma^2$ \'e
        \begin{equation}
            \hat{\sigma}^2 = - \frac{1}{2\hat{\theta}} = m_2.
        \end{equation}
        Podemos obter a estimativa de score matching $\hat{\sigma}^2$ de
        $\hat{\theta}$ desta maneira pela mesma raz\~ao que fomos
        capazes de trabalhar com par\^ametros transformados na estimativa de
        m\'axima verossimilhan\'ca.
        
        Para Gaussianas de m\'edia zero, o segundo momento $m_2$ \'e a estimativa de m\'axima
        verossimilhan\'ca da vari\^ancia, o que mostra que as estimativas de score matching e de m\'axima
        verossimilhan\'ca s\~ao as mesmas aqui. Embora os dois m\'etodos geralmente
        produzam estimativas diferentes, o resultado tamb\'em vale para
        Gaussianas multivariadas onde as estimativas de score matching tamb\'em s\~ao iguais \`as estimativas de
        m\'axima verossimilhan\'ca, veja o artigo original sobre score matching de
        \citet{Hyvarinen2005c}.
        
    \end{solution}
    
\end{exenumerate}

% --------------------------------------------------

\ex{Estimativa de m\'axima verossimilhan\'ca e modelos n\~ao normalizados}
Considere o modelo de Ising para duas vari\'aveis aleat\'orias bin\'arias $(x_1,x_2)$,
\begin{align*}
    p(x_1,x_2;\theta) \propto \exp\left(\theta x_1 x_2+x_1+x_2\right), \quad \quad x_i \in \{-1,1\},
\end{align*}

\begin{exenumerate}
    \item Compute a fun\'c\~ao de parti\'c\~ao $Z(\theta)$.
    
    \begin{solution}
        A defini\'c\~ao da fun\'c\~ao de parti\'c\~ao \'e
        \begin{align}
            Z(\theta) & = \sum_{ \{-1,1\}^2} \exp\left(\theta x_1 x_2+x_1+x_2\right).
        \end{align}
        onde temos que somar sobre $(x_1,x_2) \in  \{-1,1\}^2 = \{ (-1,1),\, (1,1),\, (1,-1),\, (-1,-1)\}$. Isso resulta em
        \begin{align}
            Z(\theta) & = \exp(-\theta-1+1)+\exp(\theta+2)+\exp(-\theta+1-1)+\exp(\theta-2)\\
            & = 2\exp(-\theta)+\exp(\theta+2)+\exp(\theta-2)
        \end{align}
        
    \end{solution}
    
    \item A figura abaixo mostra o gr\'afico de $f(\theta)= \frac{\partial \log Z(\theta)}{\partial \theta}$.\vspace{1ex}
    
    Assuma que voc\^e observe tr\^es pontos de dados $(x_1, x_2)$ iguais a
    $(-1,-1)$, $(-1,1)$, e $(1,-1)$. Usando a figura, qual \'e a
    estimativa de m\'axima verossimilhan\'ca de $\theta$? Justifique sua resposta. 
    
    \begin{center}
        \includegraphics[width=0.4\textwidth]{Ising-MLE}
    \end{center}
    
    \begin{solution}
        
        Denotando o $i$-\'esimo ponto de dado observado por $(x_1^i,x_2^i)$,
        a log-verossimilhan\'ca \'e
        \begin{align}
            \ell(\theta) & = \sum_{i=1}^n \log p(x_1^i,x_2^i; \theta)
        \end{align}
        Inserir a defini\'c\~ao de $p(x_1,x_2;\theta)$ produz
        \begin{align}
            \ell(\theta)  & = \sum_{i=1}^n \left[ \theta x_1^i x_2^i+x_1^i+x_2^i \right] - n \log Z(\theta)\\
            &= \theta \sum_{i=1}^n \left[x_1^i x_2^i\right] + \sum_{i=1}^n \left[x_1^i+x_2^i \right] - n\log Z(\theta)
        \end{align}
        Sua derivada em rela\'c\~ao a $\theta$ \'e
        \begin{align}
            \frac{\partial \ell(\theta)}{\partial \theta} & =  \sum_{i=1}^n \left[x_1^i x_2^i\right] - n  \frac{\partial \log Z(\theta)}{\partial \theta}\\
            & =  \sum_{i=1}^n \left[x_1^i x_2^i\right] - n f(\theta)
        \end{align}
        Igualando a zero, obtemos
        \begin{align}
            \frac{1}{n}  \sum_{i=1}^n \left[x_1^i x_2^i\right] &= f(\theta)
        \end{align}
        
        Uma abordagem alternativa \'e come\'car com a rela\'c\~ao mais geral
        que relaciona o gradiente da fun\'c\~ao de parti\'c\~ao
        ao gradiente do modelo n\~ao normalizado logar\'itmico. Por exemplo, se
        $$ p(\x,\thetab) = \frac{\phi(\x; \thetab)}{Z(\thetab)}$$
        temos
        \begin{align}
            \ell(\thetab) &= \sum_{i=1}^n \log p(\x_i;\thetab) \\
            & = \sum_{i=1}^n \log \phi(\x_i;\thetab) -n \log Z(\thetab)
        \end{align}
        Definir a derivada como zero resulta em
        $$\frac{1}{n} \sum_{i=1}^n \nabla_{\thetab}  \log \phi(\x_i;\thetab) = \nabla_{\thetab} \log Z(\thetab)$$
        
        Em qualquer caso, a avalia\'c\~ao num\'erica de $1/n \sum_{i=1}^n x_1^i x_2^i$ resulta em
        \begin{align}
            \frac{1}{n}  \sum_{i=1}^n \left[x_1^i x_2^i\right] & = \frac{1}{3}\left(1-1-1\right)\\
            &= -\frac{1}{3}
        \end{align}
        Pelo gr\'afico, vemos que $f(\theta)$ assume o valor $-1/3$
        para $\theta = -1$, que \'e o MLE desejado. 
        
    \end{solution}
    
\end{exenumerate}

\ex{Estimativa de par\^ametros para modelos n\~ao normalizados}
Seja $p(\x; \A) \propto \exp(-\x^\top \A \x)$ um modelo estat\'istico param\'etrico
para $\x = (x_1, \ldots, x_{100})$, onde os
par\^ametros s\~ao os elementos da matriz $\A$. Assuma que $\A$ seja
sim\'etrica e definida semi-positiva, i.e.\ $\A$ satisfaz $\x^\top
\A \x \ge 0$ para todos os valores de $\x$.

\begin{exenumerate}
    \item Para $n$ pontos de dados iid $\x_1, \ldots, \x_n$, um amigo
    prop\~oe estimar $\A$ maximizando $J(\A)$,
    \begin{equation}
        J(\A) = \prod_{k=1}^n \exp\left( - \x_k^\top \A \x_k \right).
    \end{equation}
    Explique por que este procedimento n\~ao pode fornecer estimativas de par\^ametros razo\'aveis. 
    
    \begin{solution}
        Temos que $\x_k^\top \A \x_k \ge 0$, de modo que $\exp\left( -
        \x_k^\top \A \x_k \right) \le 1$. Assim, $\exp\left( - \x_k^\top \A
        \x_k \right)$ \'e m\'aximo se os elementos de $\A$ forem zero. Isso
        significa que $J(\A)$ \'e m\'aximo se $\A = 0$, independentemente dos dados
        observados, o que n\~ao corresponde a um procedimento de estimativa
        (estimador) significativo.
        
    \end{solution}
    
    \item Explique por que a estimativa de m\'axima verossimilhan\'ca \'e f\'acil quando os $x_i$
    s\~ao n\'umeros reais, i.e. $x_i \in \mathbb{R}$, enquanto tipicamente \'e muito
    dif\'icil quando os $x_i$ s\~ao bin\'arios, i.e.\ $x_i \in \{0,1\}$.
    
    \begin{solution}
        
        Para a estimativa de m\'axima verossimilhan\'ca, precisamos normalizar o
        modelo calculando a fun\'c\~ao de parti\'c\~ao $Z(\thetab)$, que \'e
        definida como a soma/integral de $\exp(-\x^\top \A \x)$ sobre o
        dom\'inio de $\x$.
        
        Quando os $x_i$ s\~ao n\'umeros, podemos aqui obter uma
        express\~ao anal\'itica para $Z(\thetab)$. No entanto, se os $x_i$ forem bin\'arios, nenhuma
        express\~ao anal\'itica desse tipo est\'a dispon\'ivel e computar $Z(\thetab)$
        torna-se muito caro.
        
    \end{solution}
    
    \item Podemos usar score matching em vez de estimativa de m\'axima
    verossimilhan\'ca para aprender $\A$ se os $x_i$ forem bin\'arios?
    
    \begin{solution}
        N\~ao, o score matching n\~ao pode ser usado para dados bin\'arios.
    \end{solution}
    
    
\end{exenumerate}